% !TeX root = ../thesis.tex
\chapter{Erklärung zur Selbstständigkeit und zum Einsatz von \texorpdfstring{\acs{ki}}{KI}-Hilfsmitteln}
\label{chap:erklaerung}

\section*{Eigenständigkeitserklärung}

Ich erkläre hiermit, dass ich die vorliegende \thesistype\ mit dem Titel

\begin{quote}
  \emph{\thesismaintitleDE: \thesissubtitleDE}
\end{quote}

selbstständig verfasst, noch nicht anderweitig für Prüfungszwecke vorgelegt, keine anderen als die angegebenen Quellen oder Hilfsmittel benutzt sowie wörtliche und sinngemäße Zitate als solche gekennzeichnet habe.

\section*{Erklärung zum Einsatz von \texorpdfstring{\acs{ki}}{KI}-gestützten Werkzeugen}

Im Rahmen der Erstellung dieser Bachelorarbeit wurden generative Systeme der \acl{ki} (\acs{ki}) sowie \acs{ki}-gestützte Entwicklungswerkzeuge (insbesondere große Sprachmodelle und agentische Programmierassistenten) als unterstützende Hilfsmittel eingesetzt. Der Einsatz erstreckte sich auf folgende Arbeitsbereiche:

\begin{itemize}
  \item \textbf{Ideenfindung, Brainstorming und Konzeption:} Unterstützung bei der Strukturierung der Problemstellung, der Diskussion methodischer Lösungsansätze, dem Entwurf der System- und Pipeline-Architektur, der Formulierung von Fehlerhypothesen sowie der Gliederung der schriftlichen Ausarbeitung.
  \item \textbf{Softwareentwicklung, Coding, Debugging und Testautomatisierung:} Unterstützung bei der Implementierung, dem Refactoring und der Optimierung des Python-Evaluationsframeworks (\texttt{software/}), einschließlich der Vorverarbeitungs-, \acs{ocr}-, Struktur-, Prompting-, Kalibrierungs-, Matching-, Metrik- und Dashboard-Module, der Erstellung automatisierter Unit- und Integrationstests (\texttt{pytest}) sowie der technischen Fehlerdiagnose.
  \item \textbf{Aufbereitung von Auswertungen, Dokumentation und \LaTeX-Satz:} Unterstützung bei der Erstellung von Diagnose- und Auswertungsskripten, der technischen Projektdokumentation sowie der Formatierung von Tabellen, mathematischen Formeln und \LaTeX-/BibTeX-Strukturen.
  \item \textbf{Texterstellung, Formulierung und sprachliche Überarbeitung:} Unterstützung beim Verfassen und Überarbeiten von Textentwürfen, bei der sprachlichen, stilistischen und grammatikalischen Präzisierung des wissenschaftlichen Ausdrucks sowie bei der Erstellung und Überprüfung der englischsprachigen Kurzfassung (\emph{Abstract}).
\end{itemize}

Für sämtliche Einsatzbereiche gilt uneingeschränkt:
\begin{enumerate}
  \item \textbf{Eigenständige Prüfung und wissenschaftliche Verantwortung:} Alle durch \acs{ki}-Werkzeuge generierten oder überarbeiteten Vorschläge (Programmcode, Architekturentwürfe, Formeln und Textpassagen) wurden vom Verfasser eigenständig kritisch geprüft, technisch getestet, inhaltlich validiert und bei Bedarf angepasst. Die vollständige inhaltliche, methodische und wissenschaftliche Verantwortung für die Arbeit liegt beim Verfasser.
  \item \textbf{Authentizität der empirischen Ergebnisse und Quellen:} Es wurden zu keinem Zeitpunkt experimentelle Messwerte, Benchmark-Kennzahlen oder Literaturquellen durch \acs{ki} erfunden. Sämtliche in der Arbeit berichteten quantitativen Ergebnisse entstammen real ausgeführten und im Repository reproduzierbar protokollierten Läufen; alle zitierten Literaturquellen wurden anhand der Originalpublikationen verifiziert.
  \item \textbf{Wahrung von Datenschutz und Geheimhaltung:} Bei der Nutzung externer \acs{ki}-gestützter Hilfsmittel wurden die internen Datenschutz- und Vertraulichkeitsrichtlinien der BMW Group strikt eingehalten. Es wurden keinerlei personenbezogene Teilnehmerdaten, Klarnamen, Personalnummern oder ungeschwärzte interne Originaldokumente an externe Dienste übermittelt.
\end{enumerate}

\vspace{2cm}

\noindent
Deggendorf, den \submissiondate

\vspace{1.8cm}

\noindent
\rule{7cm}{0.4pt}\\
\student
